% =====================================================================
%  ĐỀ CƯƠNG ĐỒ ÁN KHOA HỌC
%  Nghiên cứu và ứng dụng mô hình Dual Agreement Consistency Learning
%  cho bài toán phân đoạn tế bào máu trong ảnh hiển vi
%  Biên dịch: pdfLaTeX (Overleaf: Menu -> Compiler -> pdfLaTeX)
% =====================================================================
\documentclass[12pt,a4paper]{article}

% ---------- Tiếng Việt ----------
\usepackage[T5]{fontenc}
\usepackage[utf8]{inputenc}
\usepackage[vietnamese]{babel}
\usepackage{lmodern}

% ---------- Trình bày ----------
\usepackage[left=3cm,right=2cm,top=2.5cm,bottom=2.5cm]{geometry}
\usepackage{setspace}\onehalfspacing
\usepackage{indentfirst}
\setlength{\parindent}{1cm}
\setlength{\parskip}{4pt}
\usepackage{titlesec}
\usepackage{enumitem}
\setlist{nosep,leftmargin=*}

% ---------- Toán, bảng, hình ----------
\usepackage{amsmath,amssymb,bm}
\usepackage{booktabs,tabularx,multirow}
\usepackage{graphicx}
\usepackage{xcolor}
\usepackage{tikz}
\usetikzlibrary{positioning,arrows.meta,fit,calc}
\usepackage{caption}
\captionsetup{font=small,labelfont=bf}
\usepackage[hidelinks]{hyperref}

\newcommand{\Lsup}{\mathcal{L}_{\mathrm{sup}}}
\newcommand{\Lcps}{\mathcal{L}_{\mathrm{cps}}}
\newcommand{\Lict}{\mathcal{L}_{\mathrm{ict}}}
\newcommand{\Ldac}{\mathcal{L}_{\mathrm{dac}}}
\newcommand{\Lkl}{\mathcal{L}_{\mathrm{KL}}}
\newcommand{\Lent}{\mathcal{L}_{\mathrm{ent}}}

% =====================================================================
\begin{document}

% ---------------------------- TRANG BÌA -----------------------------
\begin{titlepage}
  \centering
  {\large \textbf{TRƯỜNG ĐẠI HỌC \underline{\hspace{6cm}}}}\\[2pt]
  {\large \textbf{KHOA \underline{\hspace{6cm}}}}\\[2.5cm]

  {\Large \textbf{ĐỀ CƯƠNG ĐỒ ÁN KHOA HỌC}}\\[1.5cm]

  {\LARGE \textbf{NGHIÊN CỨU VÀ ỨNG DỤNG MÔ HÌNH\\[4pt]
  ``DUAL AGREEMENT CONSISTENCY LEARNING''\\[4pt]
  CHO BÀI TOÁN PHÂN ĐOẠN TẾ BÀO MÁU\\[4pt]
  TRONG ẢNH HIỂN VI}}\\[3cm]

  \begin{tabular}{ll}
    Sinh viên thực hiện: & \underline{\hspace{6cm}} \\[4pt]
    Mã số sinh viên:     & \underline{\hspace{6cm}} \\[4pt]
    Lớp:                 & \underline{\hspace{6cm}} \\[4pt]
    Giảng viên hướng dẫn:& \underline{\hspace{6cm}} \\
  \end{tabular}

  \vfill
  {\large \underline{\hspace{2.5cm}}, tháng \underline{\hspace{1cm}} năm 2026}
\end{titlepage}

\tableofcontents
\newpage

% =====================================================================
\section{Tính cấp thiết của đề tài}

Phân tích ảnh hiển vi tiêu bản máu ngoại vi là xét nghiệm cơ bản trong huyết học,
phục vụ đếm tế bào, phát hiện bất thường hình thái hồng cầu, phân loại bạch cầu và
sàng lọc các bệnh như thiếu máu, sốt rét hay bạch cầu cấp. Bước \emph{phân đoạn}
(segmentation) -- tách chính xác vùng tế bào khỏi nền -- là tiền đề cho mọi phân tích
định lượng phía sau. Khi thực hiện thủ công, công việc này tốn thời gian và phụ thuộc
vào kinh nghiệm của kỹ thuật viên.

Các mô hình học sâu như U-Net~\cite{ronneberger2015unet} đã đạt kết quả tốt trên bài
toán này~\cite{depto2021blood}, nhưng đòi hỏi lượng lớn mặt nạ (mask) được gán nhãn ở mức
điểm ảnh, vốn đắt đỏ và cần chuyên gia. Học bán giám sát (Semi-Supervised Learning, SSL)
cho phép khai thác thêm ảnh \emph{không nhãn} -- nguồn dữ liệu rất dồi dào trong phòng
xét nghiệm -- để giảm chi phí gán nhãn.

Mô hình \textbf{Dual Agreement Consistency Learning (DACL)}~\cite{wang2026dacl} là một
khung SSL mới (2026), huấn luyện đồng thời một mạng CNN siêu nhẹ (UNeXt, 1{,}47 triệu
tham số) và một mạng Transformer (Swin-Unet), kết hợp giám sát giả chéo, nhất quán nội
suy và \emph{hàm mất mát đồng thuận kép}. Trên ảnh siêu âm thai, DACL vượt các phương
pháp SSL gần đây khi chỉ có 5\% dữ liệu có nhãn. Tuy nhiên, DACL chưa được kiểm chứng
trên ảnh hiển vi -- miền dữ liệu có đặc trưng rất khác: ảnh màu nhuộm, nhiều đối tượng
nhỏ trong một ảnh, tế bào chồng lấn và chạm nhau. Đề tài này nhằm lấp khoảng trống đó.

% =====================================================================
\section{Mục tiêu nghiên cứu}

\subsection{Mục tiêu tổng quát}
Nghiên cứu, cài đặt và điều chỉnh mô hình DACL cho bài toán phân đoạn tế bào máu trong
ảnh hiển vi trong điều kiện khan hiếm nhãn, hướng tới một mô hình nhẹ có thể triển khai
trên thiết bị phổ thông.

\subsection{Mục tiêu cụ thể}
\begin{enumerate}
  \item Tìm hiểu cơ sở lý thuyết của học bán giám sát cho phân đoạn ảnh y tế: Mean Teacher,
        Cross Pseudo Supervision (CPS), Interpolation Consistency Training (ICT) và DACL.
  \item Cài đặt lại DACL bằng PyTorch và áp dụng trên bộ dữ liệu
        \emph{Blood Cell Segmentation}~\cite{depto2021blood}.
  \item Đánh giá DACL với 5\%, 10\%, 20\% dữ liệu có nhãn; so sánh với học có giám sát
        và các phương pháp SSL tiêu biểu.
  \item Thực hiện nghiên cứu cắt bỏ (ablation) từng thành phần mất mát; phân tích và đề
        xuất cải tiến thành phần đồng thuận kép.
  \item Xây dựng ứng dụng minh họa: tải ảnh lên, trả về mặt nạ phân đoạn và số lượng tế bào.
\end{enumerate}

% =====================================================================
\section{Đối tượng và phạm vi nghiên cứu}

\begin{itemize}
  \item \textbf{Đối tượng:} các phương pháp học bán giám sát cho phân đoạn ảnh, trọng tâm
        là DACL; ảnh hiển vi quang học tiêu bản máu ngoại vi nhuộm màu.
  \item \textbf{Phạm vi:} phân đoạn ngữ nghĩa 2D (tế bào/nền); tách và đếm tế bào bằng hậu
        xử lý. Không bao gồm phân loại chi tiết các dòng bạch cầu và không triển khai lâm sàng.
\end{itemize}

% =====================================================================
\section{Tổng quan về mô hình DACL}

\subsection{Kiến trúc}
DACL gồm ba mạng (Hình~\ref{fig:dacl}):
\begin{itemize}
  \item $f^1_\theta$ -- \textbf{UNeXt}~\cite{valanarasu2022unext}: CNN kết hợp MLP, số kênh
        bộ mã hóa $\{32,64,128,160,256\}$, 1{,}47M tham số, 7{,}03 GFLOPs. Đây là mô hình
        \emph{duy nhất được dùng khi suy luận}.
  \item $f^2_\phi$ -- \textbf{Swin-Unet tiny}~\cite{cao2022swinunet}: Transformer, 27{,}15M
        tham số, khởi tạo từ trọng số tiền huấn luyện, cửa sổ dịch chuyển kích thước 7.
  \item $f^{\mathrm{ema}}_{\bar\phi}$ -- \textbf{giáo viên EMA}: trung bình trượt mũ của
        tham số $\phi$, theo Mean Teacher~\cite{tarvainen2017mt}.
\end{itemize}

\begin{figure}[ht]
\centering
\begin{tikzpicture}[
  node distance=8mm and 14mm,
  data/.style={draw,rounded corners,fill=gray!12,minimum width=24mm,minimum height=8mm,align=center,font=\small},
  net/.style={draw,thick,rounded corners,minimum width=26mm,minimum height=9mm,align=center,font=\small},
  loss/.style={draw,dashed,red!70!black,rounded corners,align=center,font=\small,inner sep=4pt},
  arr/.style={-{Stealth[length=2mm]},thick}
]
  \node[data] (xl) {Ảnh có nhãn\\$X_L, Y_{gt}$};
  \node[data,below=14mm of xl] (xu) {Ảnh không nhãn\\$X_{UL}$};
  \node[data,below=10mm of xu] (xm) {Mixup\\$\hat X_{UL}$};

  \node[net,fill=blue!10,right=of xl] (cnn) {UNeXt $f^1_\theta$\\(CNN nhẹ)};
  \node[net,fill=orange!15,right=of xu] (tr) {Swin-Unet $f^2_\phi$\\(Transformer)};
  \node[net,fill=orange!5,right=of xm] (ema) {Giáo viên EMA\\$f^{\mathrm{ema}}_{\bar\phi}$};

  \draw[arr] (xl) -- (cnn);
  \draw[arr] (xu) -- (tr);
  \draw[arr] (xu.north east) -- (cnn.south west);
  \draw[arr] (xu) -- (xm);
  \draw[arr] (xm.north east) -- (tr.south west);
  \draw[arr] (xu.south east) -- (ema.north west);
  \draw[arr,dotted] (tr) -- node[right,font=\scriptsize]{EMA} (ema);

  \node[loss,right=16mm of cnn] (lsup) {$\Lsup$\\(CE + Dice)};
  \node[loss,right=16mm of tr]  (lcps) {$\Lcps$ \;\&\; $\Ldac$};
  \node[loss,right=16mm of ema] (lict) {$\Lict$\\(MSE)};

  \draw[arr] (cnn) -- (lsup);
  \draw[arr] (cnn.south east) -- (lcps.north west);
  \draw[arr] (tr) -- (lcps);
  \draw[arr] (ema) -- (lict);
  \draw[arr] (tr.south east) -- (lict.north west);
\end{tikzpicture}
\caption{Sơ đồ tổng quát khung DACL (vẽ lại theo~\cite{wang2026dacl}). Đường nét đứt màu
đỏ là các hàm mất mát.}
\label{fig:dacl}
\end{figure}

\subsection{Giám sát giả chéo (CPS)}
Với ảnh không nhãn $X_{UL}$, mỗi mạng sinh nhãn giả one-hot
$\tilde Y^1=f_{OH}\big(f^1_\theta(X_{UL})\big)$, $\tilde Y^2=f_{OH}\big(f^2_\phi(X_{UL})\big)$
(không lan truyền gradient qua nhãn giả) để giám sát mạng còn lại~\cite{chen2021cps}:
\begin{align}
  \mathcal{L}^1_{\mathrm{cps}} &= H\big(f^1_\theta(X_{UL}),\tilde Y^2\big) + D\big(f^1_\theta(X_{UL}),\tilde Y^2\big),\\
  \mathcal{L}^2_{\mathrm{cps}} &= H\big(f^2_\phi(X_{UL}),\tilde Y^1\big) + D\big(f^2_\phi(X_{UL}),\tilde Y^1\big),
\end{align}
trong đó $H$ là cross-entropy và $D$ là Dice loss.

\subsection{Nhất quán nội suy (ICT)}
Với $\hat X_{UL}=\mathrm{Mix}(X^i_{UL},X^j_{UL})$ (hệ số trộn $\sigma=0{,}5$), mạng
Transformer được ép khớp với phép trộn tương ứng của các dự đoán từ giáo viên
EMA~\cite{verma2022ict}:
\begin{equation}
  \Lict=\sum\Big(f^2_\phi(\hat X_{UL})-\mathrm{Mix}\big(f^{\mathrm{ema}}_{\bar\phi}(X^i_{UL}),
  f^{\mathrm{ema}}_{\bar\phi}(X^j_{UL})\big)\Big)^2 .
\end{equation}

\subsection{Đồng thuận kép (DAC) -- đóng góp chính của DACL}
Gọi $p(u)$, $q(u)$ là phân phối lớp tại điểm ảnh $u$ do $f^1_\theta$ và $f^2_\phi$ dự đoán:
\begin{align}
  \Lkl  &= \mathbb{E}\Big[\tfrac{1}{HW}\textstyle\sum_u \mathrm{KL}\big(p(u)\,\|\,q(u)\big)\Big],\\
  \Lent &= \mathbb{E}\Big[\tfrac{1}{HW}\textstyle\sum_u
           H\big(p(u)\big)+H\big(q(u)\big)-H\big(\tfrac{p(u)+q(u)}{2}\big)\Big],\\
  \Ldac &= \Lkl + \Lent .
\end{align}

\paragraph{Nhận xét phân tích.} Vì độ phân kỳ Jensen--Shannon
$\mathrm{JSD}(p,q)=H\big(\tfrac{p+q}{2}\big)-\tfrac12\big[H(p)+H(q)\big]$, ta có
\begin{equation}
  H(p)+H(q)-H\big(\tfrac{p+q}{2}\big)=\tfrac12\big[H(p)+H(q)\big]-\mathrm{JSD}(p,q).
\end{equation}
Như vậy $\Lent$ vừa khuyến khích dự đoán \emph{tự tin} (entropy thấp), vừa có xu hướng
\emph{tăng} độ khác biệt JSD giữa hai mạng; số hạng $\Lkl$ đóng vai trò đối trọng. Đề tài
sẽ khảo sát thực nghiệm các biến thể thay thế (mục~\ref{sec:improve}).

\subsection{Hàm mục tiêu tổng}
\begin{equation}
  \mathcal{L}_{\mathrm{total}}=\big(\mathcal{L}^1_{\mathrm{sup}}+\mathcal{L}^2_{\mathrm{sup}}\big)
  +\lambda\big(\mathcal{L}^1_{\mathrm{cps}}+\mathcal{L}^2_{\mathrm{cps}}\big)
  +\tau\,\Lict+\beta\,\Ldac ,
\end{equation}
với $\{\lambda,\tau,\beta\}=\{2{,}0;\,2{,}0;\,10{,}0\}$ theo bài báo gốc;
$\mathcal{L}_{\mathrm{sup}}$ là tổng CE và Dice trên dữ liệu có nhãn.

% =====================================================================
\section{Dữ liệu}

\subsection{Bộ dữ liệu Blood Cell Segmentation}
Đề tài sử dụng bộ dữ liệu công bố kèm công trình của Depto và cộng sự~\cite{depto2021blood}
(kho mã: \url{https://github.com/Deponker/Blood-cell-segmentation}), được giới thiệu là
bộ dữ liệu phân đoạn tế bào máu hiển vi lớn nhất tại thời điểm công bố:
\begin{itemize}
  \item 1\,328 ảnh hiển vi tiêu bản máu gốc và 1\,328 mặt nạ ground truth tương ứng
        (tổng 2\,656 tệp);
  \item kho mã được phát hành theo giấy phép MIT; dữ liệu tải qua liên kết Google Drive
        trong README;
  \item độ phân giải ảnh, định dạng tệp và cách mã hóa nhãn trong mặt nạ sẽ được kiểm tra
        bằng script \texttt{scripts/inspect\_dataset.py} sau khi tải về
        (\textit{dự kiến: mặt nạ nhị phân tế bào/nền}).
\end{itemize}

\subsection{Phân chia dữ liệu}
Dữ liệu được chia ngẫu nhiên, cố định hạt giống (seed) theo tỉ lệ 70/10/20
(Bảng~\ref{tab:split}). Từ tập huấn luyện, chỉ giữ nhãn cho 5\%, 10\%, 20\% ảnh; phần
còn lại được coi là ảnh không nhãn.

\begin{table}[ht]
\centering
\caption{Phân chia dữ liệu dự kiến.}
\label{tab:split}
\begin{tabular}{lcccc}
\toprule
Tập & Số ảnh & 5\% nhãn & 10\% nhãn & 20\% nhãn \\
\midrule
Huấn luyện (có nhãn / không nhãn) & 930 & 46 / 884 & 93 / 837 & 186 / 744 \\
Kiểm định (validation)            & 133 & \multicolumn{3}{c}{133} \\
Kiểm thử (test)                   & 265 & \multicolumn{3}{c}{265} \\
\midrule
Tổng                              & 1\,328 & & & \\
\bottomrule
\end{tabular}
\end{table}

\subsection{Tiền xử lý và tăng cường}
\begin{itemize}
  \item Đưa ảnh về $224\times224$ (tương thích cửa sổ 7 của Swin-Unet; có thể thử
        $448\times448$ nếu đủ tài nguyên), chuẩn hóa theo thống kê ImageNet.
  \item Tăng cường cho ảnh có nhãn: lật, xoay $90^\circ$, co giãn ngẫu nhiên, thay đổi
        độ sáng/tương phản/màu (color jitter) để mô phỏng khác biệt nhuộm.
  \item Tùy chọn: chuẩn hóa màu nhuộm (Macenko/Reinhard).
\end{itemize}

% =====================================================================
\section{Phương pháp thực hiện}

\subsection{Các điều chỉnh của DACL cho ảnh hiển vi}
\begin{enumerate}
  \item \textbf{Đầu ra:} số lớp $C$ theo nhãn thực tế của bộ dữ liệu (dự kiến $C=2$).
  \item \textbf{Tách tế bào chạm nhau:} hậu xử lý bằng biến đổi khoảng cách và
        watershed để tách từng tế bào, phục vụ đếm.
  \item \textbf{Đơn vị đo:} HD95 và ASD tính theo điểm ảnh (không có thông tin
        kích thước vật lý của điểm ảnh).
  \item \textbf{Siêu tham số:} khảo sát lại $\beta\in\{1,5,10\}$ và lịch tăng dần
        (ramp-up) trọng số các mất mát không giám sát.
\end{enumerate}

\subsection{Đề xuất cải tiến thành phần đồng thuận}\label{sec:improve}
\begin{itemize}
  \item \textbf{DAC-sym:} thay $\mathrm{KL}(p\|q)$ bằng KL đối xứng hoặc JSD.
  \item \textbf{DAC-ent:} chỉ giữ phần entropy $\tfrac12[H(p)+H(q)]$, bỏ $-\mathrm{JSD}$.
  \item \textbf{DAC-conf:} chỉ tính DAC trên các điểm ảnh có độ tin cậy
        $\max_c p_c(u) > \gamma$.
\end{itemize}

\subsection{Thiết lập thực nghiệm}
\begin{itemize}
  \item \textbf{Huấn luyện:} SGD, lr $=0{,}01$, momentum $0{,}9$, weight decay $10^{-4}$;
        batch có nhãn 1 (hoặc 2), batch không nhãn 4; 200--400 epoch; lưu trọng số UNeXt
        tốt nhất trên tập kiểm định.
  \item \textbf{Phương pháp so sánh:} UNeXt và Swin-Unet học có giám sát (giới hạn
        dưới/trên), Mean Teacher~\cite{tarvainen2017mt}, ICT~\cite{verma2022ict},
        CPS~\cite{chen2021cps}, CTCT~\cite{luo2022ctct}, DACL và các biến thể đề xuất.
  \item \textbf{Ablation:} Baseline $\rightarrow$ +$\Lcps$ $\rightarrow$
        +$\Lcps$+$\Lict$ $\rightarrow$ +$\Lcps$+$\Ldac$ $\rightarrow$ đầy đủ.
  \item \textbf{Chỉ số:} DSC, Jaccard (IoU), HD95, ASD; số tham số, GFLOPs, tốc độ
        suy luận (FPS); sai số đếm tế bào (MAE). Kiểm định Wilcoxon cho ý nghĩa thống kê.
\end{itemize}

\subsection{Công cụ}
Python, PyTorch, MONAI/MedPy (tính chỉ số), Albumentations; mã nguồn tham khảo SSL4MIS,
UNeXt và Swin-Unet; quản lý mã nguồn bằng Git/GitHub; huấn luyện trên GPU cá nhân hoặc
Google Colab/Kaggle.

% =====================================================================
\section{Kết quả dự kiến}
\begin{itemize}
  \item Mã nguồn DACL cho ảnh hiển vi tế bào máu, công khai trên GitHub.
  \item Bảng so sánh định lượng và hình ảnh trực quan giữa các phương pháp ở các mức nhãn.
  \item Mô hình UNeXt ($\approx$1{,}5M tham số) đạt chất lượng tiệm cận học có giám sát
        đầy đủ khi chỉ dùng 10--20\% nhãn.
  \item Ứng dụng minh họa phân đoạn và đếm tế bào.
  \item Báo cáo đồ án và (nếu kết quả tốt) bài báo hội nghị trong nước.
\end{itemize}

% =====================================================================
\section{Kế hoạch thực hiện}

\begin{table}[ht]
\centering
\caption{Tiến độ dự kiến (16 tuần).}
\begin{tabularx}{\textwidth}{c X}
\toprule
Tuần & Nội dung công việc \\
\midrule
1--2   & Nghiên cứu tài liệu; tải, kiểm tra và phân chia bộ dữ liệu \\
3--5   & Cài đặt và huấn luyện UNeXt, Swin-Unet có giám sát \\
6--9   & Cài đặt DACL và các phương pháp SSL so sánh \\
10--12 & Thực nghiệm ở các mức nhãn; ablation; thử các biến thể cải tiến \\
13--14 & Hậu xử lý tách/đếm tế bào; xây dựng ứng dụng minh họa \\
15--16 & Viết báo cáo, chuẩn bị bảo vệ \\
\bottomrule
\end{tabularx}
\end{table}

% =====================================================================
\section{Bố cục dự kiến của báo cáo}
\begin{enumerate}
  \item[] \textbf{Mở đầu}
  \item[] \textbf{Chương 1.} Tổng quan về ảnh hiển vi tế bào máu và bài toán phân đoạn
  \item[] \textbf{Chương 2.} Cơ sở lý thuyết: mạng phân đoạn và học bán giám sát
  \item[] \textbf{Chương 3.} Mô hình DACL và các điều chỉnh đề xuất
  \item[] \textbf{Chương 4.} Thực nghiệm và đánh giá
  \item[] \textbf{Chương 5.} Ứng dụng minh họa
  \item[] \textbf{Kết luận và hướng phát triển}
\end{enumerate}

% =====================================================================
\bibliographystyle{unsrt}
\bibliography{references}

\end{document}
